\writeSectionFrame{\presentationTitle}{Übung}
\begin{frame}
    \frametitle{Ziel der Übung}
    \begin{itemize}
        \item Implementieren Sie das Strategie-Muster, um verschiedene Angriffsstrategien in einem Kampfspiel zu modellieren.
        \item Spieler und Gegner wechseln ihre Angriffsstrategie zufällig bei jedem Zug.
        \item Der Kampf dauert an, bis einer der Kämpfer besiegt ist.
    \end{itemize}
\end{frame}

\begin{frame}
    \frametitle{Anforderungen}
    \begin{itemize}
        \item Modellieren Sie drei Angriffsstrategien:
        \begin{itemize}
            \item Nahkampf (Melee)
            \item Fernkampf (Ranged)
            \item Magie (Magic)
        \end{itemize}
        \item Verwenden Sie eine Schnittstelle \texttt{AttackStrategy}, die eine Methode \texttt{attack(Player target)} deklariert.
        \item Jeder Spieler hat Lebenspunkte (\texttt{health}), die durch Angriffe verringert werden.
        \item Legen Sie am einfachsten für jede Angriffsart eine Menge Lebenspunkten fest, die angezogen werden
        \item Sobald ein Spieler besiegt ist (Lebenspunkte <= 0), soll er eine \texttt{DeadStrategy} verwenden.
    \end{itemize}
\end{frame}

\begin{frame}
    \frametitle{Schritte zur Lösung}
    \begin{itemize}
        \item Definieren Sie die Schnittstelle \texttt{AttackStrategy} und implementieren Sie die konkreten Strategien \texttt{MeleeStrategy}, \texttt{RangedStrategy} und \texttt{MagicStrategy}.
        \item Erstellen Sie eine Klasse \texttt{Player} mit:
        \begin{itemize}
            \item Name und Lebenspunkten.
            \item Eine Methode \texttt{takeDamage(int amount)}, um den Schaden zu verarbeiten.
            \item Eine Methode \texttt{performAttack(Player target)}, die die aktuelle Strategie verwendet.
        \end{itemize}
        \item Implementieren Sie die Klasse \texttt{GameContext}, die den Spielverlauf steuert:
        \begin{itemize}
            \item Eine Methode \texttt{switchStrategy(Player player)}, die eine zufällige Strategie wählt.
            \item Eine Methode \texttt{startGame()}, die den Kampf bis zum Ende durchführt.
        \end{itemize}
    \end{itemize}
\end{frame}

\begin{frame}
    \frametitle{Zusatzanforderungen}
    \begin{itemize}
        \item Implementieren Sie eine \texttt{DeadStrategy}, die signalisiert, dass ein Spieler besiegt ist und nicht mehr kämpfen kann.
        \item Verwenden Sie die Klasse \texttt{Random}, um zufällig eine Strategie auszuwählen.
        \item Lassen Sie das Spiel automatisch bis zum Ende ablaufen.
    \end{itemize}
\end{frame}

\begin{frame}
    \frametitle{Erweiterungsmöglichkeiten}
    \begin{itemize}
        \item Fügen Sie zusätzliche Strategien oder Spezialfähigkeiten hinzu, um den Kampf spannender zu gestalten.
        \item Ermöglichen Sie es dem Spieler, seine Strategie manuell zu wählen.
        \item Implementieren Sie ein einfaches Text-Interface, um den Kampfverlauf darzustellen.
    \end{itemize}
\end{frame}

\begin{frame}
    \frametitle{Hinweise}
    \begin{itemize}
        \item Achten Sie auf die korrekte Anwendung des Strategie-Musters: Die konkrete Strategie sollte unabhängig von der \texttt{Player}-Klasse implementiert sein.
        \item Überlegen Sie sich, wie das \texttt{GameContext}-Objekt den Ablauf des Spiels steuern kann.
        \item Testen Sie das Spiel mit verschiedenen Lebenspunkten und Strategiekombinationen.
    \end{itemize}
\end{frame}

\begin{frame}
    \frametitle{Beispielhafter Ablauf}
Spieler 1 erleidet 15 Schaden. Verbleibende HP: 35

Spieler 1 wählt eine neue Strategie.

Schwingt das Schwert!

Gegner erleidet 10 Schaden. Verbleibende HP: 25

Gegner wählt eine neue Strategie.

Wirkt einen Feuerball!

Spieler 1 erleidet 15 Schaden. Verbleibende HP: 20

Spieler 1 wählt eine neue Strategie.

Wirkt einen Feuerball!

Gegner erleidet 15 Schaden. Verbleibende HP: 10

Gegner wählt eine neue Strategie.

Schwingt das Schwert!

Spieler 1 erleidet 10 Schaden. Verbleibende HP: 10
\end{frame}

\begin{frame}
    \frametitle{Beispielhafter Ablauf}
Spieler 1 wählt eine neue Strategie.

Schießt einen Pfeil!

Gegner erleidet 7 Schaden. Verbleibende HP: 3

Gegner wählt eine neue Strategie.

Schwingt das Schwert!

Spieler 1 erleidet 10 Schaden. Verbleibende HP: 0

Spieler 1 ist besiegt und kann nicht mehr angreifen.

Gegner gewinnt!
\end{frame}